\section{Q\&A 延伸：适用范围、数学假设与未完成比较}

Q\&A 没有新增独立 slide，却补齐了正文最容易被夸大的四个边界：warm-start evidence 的 architecture、非 softmax layer 的覆盖、non-negative assumption，以及 packing comparison 的缺失。

\subsection{Mixed modalities 可以想象，但每种 layer 都要满足条件}

若 image features 与 text features 都经过非负处理，并进入 softmax layer，同一 generalized co-occurrence 形式在数学上可以累计 mixed inputs。可写成
\[
H_m=\operatorname{concat}(H_m^{\text{image}},H_m^{\text{text}}),
\qquad F=H^\top Y.
\]
但真实 multimodal model 还含 convolution/vision encoder、normalization、projection、cross-attention 等组件，不能只初始化最后一层就声称完整模型被显式优化。

\teachervoice{01:14:33--01:15:25，讲者认为 mixed non-negative features 可共享同一 softmax warm start，因此 image captioning 在概念上可能；这只是可推导方向，不是已完成系统。}

\subsection{Non-negativity 为什么不是实现细节}

公式含 $\log F_{j,i}$。若 feature 有负值，outer-product sums 可能为负，log 不再对应概率或共现强度。可以通过 feature splitting、positive transform 或另行推导处理，但每种选择都会改变统计意义和最优条件。

\begin{importantbox}{非负假设的直接检查}
在实现显式初始化前，先检查 feature minimum、negative fraction、zero columns 和 column mass。不要在出现负数后静默取 absolute value；那会改变方向信息且不再对应原推导。
\end{importantbox}

\teachervoice{01:15:29--01:16:14，Williams 用“negative value 进入 logarithm 会发生什么”回答假设来源，并明确说负 feature 需要新的 derivation，而不是小修补。}

\subsection{标准 attention、SAFFU 与 concatenation 的适用场景}

三种机制解决不同问题。Concatenation 保留每个位置的独立 feature slot，几乎不需重新加权，但维度随 context 线性增长；standard attention 把不同位置 superpose 后再学习选择；SAFFU 直接变换 comparison matrix，并围绕该矩阵做显式初始化与缓存。

\begin{knowledgebox}{三种 context 表示对照}
\begin{tabularx}{\textwidth}{L{0.19\textwidth} L{0.25\textwidth} X}
\toprule
方式 & 优点 & 主要代价/限制 \\
\midrule
Concatenation & 位置信息完全分离、易做单层实验 & 维度随 context 增长，不适合长序列 \\
Standard attention & 灵活 projection、多头、适配不同表示 & 随机参数多，quadratic comparisons \\
SAFFU proposal & 可显式初始化、可缓存固定比较 & 强假设、非标准实现、composition 未完全解决 \\
\bottomrule
\end{tabularx}
\end{knowledgebox}

\subsection{本章小结}

Q\&A 把课程从“惊人的效率故事”拉回可验证边界：实验不全是 attention，跨 layer type 需要新推导，negative features 破坏当前 log 公式，packing 比较尚未完成，软件也仍在建设。

